% sections/abstract.tex  (input in the preamble, before \begin{document})
\abstract{%
The items of a psychometric scale constitute a proposal, written in ordinary
language, about the structure of a construct---and language-model embeddings
make that proposal measurable before a single respondent is recruited. We
introduce \texttt{semanticfa}, an R package for semantic factor analysis
(SFA): exploratory factor analysis of the similarity structure of
language-model embeddings of scale items, requiring no response data at any
stage. The package implements the complete workflow: similarity construction
under keyed, keying-free, and natural-language-inference encodings;
factor retention by multi-criterion consensus; extraction with Monte
Carlo--calibrated diagnostics; interpretation by anchoring, semantic
projection, congruence, and jingle--jangle screening; and response-free
refinement (redundancy detection, short forms, candidate-item vetting). A
fully reproducible demonstration
exercises every exported function on the 50-item International Personality
Item Pool Big-Five Factor Markers, embedding the items with
Qwen3-Embedding-8B and validating the semantic structure against a
conventional factor analysis of \valNHuman{} respondents. The retention consensus
indicated \valConsensusKWord{} factors, and the semantic solution recovered the Big Five
domains, with congruence to the human factors graded by domain (Tucker's
$\phi$ from \valPhiO{} for Openness to \valPhiA{} for Agreeableness) and semantic
similarities correlating $r = \valPairR$ with response correlations across all
\valNPairs{} item pairs. Sign-flipping reverse-keyed embeddings produced anti-topic
vectors that aligned worst with keyed human data, and a semantically central
25-item short form recovered the theoretical partition better than the full
inventory, though parallel analysis then retained \valSfKRedWord{} factors rather than
\valSfKFullWord. We frame the agreement of the two structures---semantic--behavioral
coherence---as an estimable relation in a construct's nomological network: a
complement to, not a replacement for, response-based factor analysis. The
package is available on CRAN.
}
\keywords{semantic factor analysis, language models, embeddings, construct validity, scale development, R package}
